\section{The integrability of curvature equations and complex ODEs}

\label{sec-MP-E}

In this section, we briefly review how to connect the integrability of the curvature equation~\eqref{00eq: DE on C} with a class of second order ODEs in complex variables; see, e.g., \cite{Chai-Lin-Wang,Coaxial}. Given a solution $u(x)$ of~\eqref{00eq: DE on C}, we define
\begin{equation*}%\label{e22q-3}
 Q(x):=-\frac12\biggl(u_{xx}(x)-\frac12 u_x(x)^2\biggr),\qquad
 x\in\mathbb{C}.
\end{equation*}
Then it is easy to see $Q_{\bar{x}}(x)\equiv 0$, so $Q(x)$ is a meromorphic function in $\mathbb{C}$ with poles belonging to
\[
 \mathcal{I}:=\{0,1,t_1,\dots,t_{n+m}\}.
\]
More precisely, at each $p\in\mathcal{I}$, since $\Delta (\ln|x-p|)=2\pi\delta_p$, by~\eqref{00eq: DE on C} and a standard elliptic regularity argument (cf. \cite{BM}), we obtain that $u(x)-2\alpha_p\ln|x-p|$ is bounded near $p$, i.e.,
\[u(x)=2\alpha_p\ln|x-p|+O(1)\qquad\text{for $x$ near $p$}.\]
Consequently,
\[Q(x)=\frac{\alpha_p(\alpha_p+2)}{4}(x-p)^{-2}+O\bigl((x-p)^{-1}\bigr)\qquad\text{for $x$ near $p$},\]
and
\[Q(x)=\frac{\alpha_\infty(\alpha_\infty+2)}{4}x^{-2}+O\bigl(x^{-3}\bigr)\qquad\text{for $x\to \infty$}.\]
Thus,
\begin{align}\label{e0q-30}
Q(x)=\sum_{p\in\mathcal{I}}\biggl[\frac{\beta_p}{(x-p)^2}+\frac{d_p}{(x-p)}\biggr],
\end{align}
where
\begin{equation}\label{e0q-31}
 \beta_p:=\frac{\alpha_{p}(\alpha_{p}+2)}{4},\qquad p\in \{0,1,t_1,\dots, t_{n+m},\infty\},
\end{equation}
and $d_p$'s are some constants satisfying
\begin{align}\label{e0q-32}
\sum_{p\in \mathcal{I}}d_p=0,\qquad \sum_{p\in \mathcal{I}}(\beta_p+pd_p)=\beta_\infty.
\end{align}
Thus, $d_0$ and $d_1$ can be uniquely determined by
\[
 \boldsymbol d:=(d_{t_1},\dots, d_{t_{n+m}}),\qquad
 \bt:=(t_1,\dots,t_{n+m})
\]
as
\begin{align*}%\label{e0q-33}
d_1=\beta_\infty-\sum_{p\in\mathcal{I}\setminus\{0,1\}}(\beta_p+pd_p),
 \qquad d_0=\sum_{p\in \mathcal{I}\setminus\{0,1\}}
 (\beta_p+(p-1)d_p)-\beta_\infty,
\end{align*}
 while $\boldsymbol d$ could be considered as $n+m$ unknown coefficients.

On the other hand, as mentioned in the introduction, the Liouville
theorem (see, e.g., \cite{Chai-Lin-Wang}) says that for a solution
$u(x)$ of~\eqref{00eq: DE on C}, there is a developing map $h(x)$ such
that
\begin{equation}\label{eqLivoulle-c1}
u(x)=\ln\frac{8|h'(x)|^2}{(1+|h(x)|^2)^2},\qquad x\in \dot{\C}:=\mathbb{C}\setminus\mathcal{I}.
\end{equation}
The developing map $h(x)$ is multi-valued in $\mathbb{C}$ and due to
\eqref{eqLivoulle-c1}, \emph{the projective monodromy group of $h(x)$
 is contained in $\PSU(2)$}. More precisely, given any $\ell\in
\pi_1\bigl(\dot{\C}, q_0\bigr)$ where $q_0\in\dot{\C}$ is a base point, the
analytic continuation of $h(x)$, denoted by $\ell^*h(x)$, is also a
developing map of $u(x)$ and so $\ell^*h(x)=\frac{ah(x)+b}{ch(x)+d}$
for some projective monodromy matrix $\SM{a}{b}{c}{d}\in\PSU(2)$.

To connect $Q(x)$ and $h(x)$, we note that~\eqref{eqLivoulle-c1} implies
\begin{equation}\label{e22q-1}
\{h(x),x\}:=\biggl(\frac{h''(x)}{h'(x)}\biggr)'-\frac12 \biggl(\frac{h''(x)}{h'(x)}\biggr)^2=u_{xx}(x)-\frac12 u_x(x)^2=-2Q(x),
\end{equation}
where $\{h(x),x\}$ is the Schwarz derivative. We refer the reader to
\cite{Chai-Lin-Wang} for details of computation.

Consider the following second order linear ODE with the complex variable $x$:
\begin{equation}\label{00eq: ODE on C}
y''(x)=Q(x)y(x),\qquad x\in\mathbb{C}.
\end{equation}
Then a classical result (see \cite{Whittaker-Watson}) says that for
any basis $(\hat{y}_1,\hat{y}_2)$ of~\eqref{00eq: ODE on C}, the
Schwarz derivative $\{\hat{y}_2/\hat{y}_1,x\}$ always satisfies
\begin{equation*}
\{\hat{y}_2/\hat{y}_1,x\}=-2Q(x).
\end{equation*}
From this and~\eqref{e22q-1}, we conclude that if~\eqref{00eq: ODE on
 C} is derived from a solution of~\eqref{00eq: DE on C}, then there
is a basis $(y_1,y_2)$ of~\eqref{00eq: ODE on C} such that
\begin{equation}\label{e22q-2}
h(x)=y_2(x)/y_1(x).
\end{equation}
Furthermore, the projective monodromy group of~\eqref{00eq: ODE on C}
with respect to $y_1$ and $y_2$ is the same as that of $h(x)$, i.e., is
contained in $\PSU(2)$.

Conversely, given $Q(x)$ via~\eqref{e0q-30}--\eqref{e0q-32}, if the
corresponding ODE~\eqref{00eq: ODE on C} has a basis $(y_1,y_2)$ such
that the projective monodromy group of the ratio $h(x)=y_2(x)/y_1(x)$
is contained in $\PSU(2)$, then
$u(x):=\ln\frac{8|h'(x)|^2}{(1+|h(x)|^2)^2}$ is a solution of
\eqref{00eq: DE on C}.


For the reader's convenience, we now briefly recall some basic notions
of ODEs like~\eqref{00eq: ODE on C}. We refer the reader to
\cite{Yoshida,Whittaker-Watson} for a comprehensive
introduction. Generally,~\eqref{00eq: ODE on C} is called
\emph{Fuchsian} because the pole order of $Q(x)$ is at most $2$ at any
singularities.
Note that the Riemann scheme (see \cite[p.~11]{Yoshida} for the definition of the Riemann scheme) of~\eqref{00eq: ODE on C} is
 given by
\begin{equation}\label{00e0q-10}\begin{pmatrix}
0&1&t_1&\cdots& t_{n+m}&\infty\\
-\frac{{\alpha}_0}{2} &-\frac{{\alpha}_1}{2}&-\frac{{\alpha}_{t_1}}{2}&\cdots&-\frac{{\alpha}_{t_{n+m}}}{2}&-(\frac{{\alpha}_\infty}{2}+1)\\
\frac{{\alpha}_0}{2}+1&\frac{{\alpha}_1}{2}+1&\frac{{\alpha}_{t_1}}{2}+1&\cdots&\frac{{\alpha}_{t_{n+m}}}{2}+1&\frac{{\alpha}_\infty}{2}
\end{pmatrix}.\end{equation}
Observe that the exponent difference at $t_j$ is an integer for
$n+1\leq j\leq n+m$.
Since $Q(x)$ comes from a solution $u(x)$ of (\ref{00eq: DE on C}), it is well known that
(\ref{00eq: ODE on C}) is apparent at $t_j$ \big(i.e., no logarithmic
singularity near $t_j$\big) and so
the local monodromy matrix of~\eqref{00eq: ODE on C} at $t_j$ is
$(-1)^{{\alpha}_{t_j}}I_2$ for all $n+1\leq j\leq n+m$.
% See Section ? for the explanation of these facts.
In this paper, we will also consider a Fuchsian differential equation
\eqref{00eq: ODE on C} which may not come from a solution of
\eqref{00eq: DE on C}. In this case, unless otherwise specified,
\emph{we always assume that~\eqref{00eq: ODE on C} with the Riemann
 scheme~\eqref{00e0q-10} is apparent at $t_j$ for all $n+1\leq j\leq
 n+m$}.

It is well known that the necessary and sufficient condition of
\eqref{00eq: ODE on C} being apparent at $t_j$ for all $n+1\leq j\leq
n+m$ can be derived from the Frobenius method.

\begin{Theorem} \label{prop: apparent}
 There are polynomials
 \[
 \hat{\mathcal{P}}_j(\boldsymbol d,\boldsymbol t)\in\Q_\theta(\bt)[\boldsymbol d]=d_{t_{j}}^{\theta_{t_j}}+ \text{l.o.t.\ in $\boldsymbol d$} \in \Q_\theta(\boldsymbol t)[\boldsymbol d]
 \]
 of total degree $\theta_{t_j}$ in $\boldsymbol d$, $j=n+1,\dots,n+m$,
 such that~\eqref{00eq: ODE on C}
 is apparent at $t_j$ for all $n+1\leq j\leq n+m$ if
 and only if $\boldsymbol{d}$ and $\bt$ satisfy $\hat{P}_j(\boldsymbol{d},\bt)=0$ for all
 $j=n+1,\dots,n+m$.

 Furthermore,
 $\hat{\mathcal{P}}_j(\boldsymbol{d},\bt)$ can be written as
 \begin{equation}\label{e22q:PP}\hat{\mathcal{P}}_j(\boldsymbol{d},\bt)=\frac{\mathcal{P}_j(\boldsymbol{d},\bt)}
 {\prod_{p\in\mathcal{I}\setminus\{t_j\}}(p-t_j)^{N_j}}\end{equation}
 for some $N_j\in\mathbb{Z}_{\geq 0}$ and $\mathcal{P}_j(\boldsymbol{d},\bt)\in \Q_{\theta}[\boldsymbol{d},\bt]$ such that $\mathcal{P}_j(\boldsymbol{d},\bt)$ and $\prod_{p\in\mathcal{I}\setminus\{t_j\}}(p-t_j)^{N_j}$ are coprime. Thus,~\eqref{00eq: ODE on C}
 is apparent at $t_j$ for all $n+1\leq j\leq n+m$ if
 and only if $\mathcal{P}_j(\boldsymbol{d},\bt)=0$ for all
 $j=n+1,\dots,n+m$.
\end{Theorem}

\begin{proof} Let us take $t_{n+m}$ for example and in the following proof, we write $t_{n+m}=t$ for convenience. By the Frobenius method, (\ref{00eq: ODE on C}) is apparent at the singularity $t$ if and only if it has a~solution of the form
\[y(x)=\sum_{j=0}^{\infty}c_j u^{j-\alpha_{t}/2},\qquad u=x-t,\quad c_0=1.\]
Observe that
{\allowdisplaybreaks
\begin{align*}
Q(x)&=\sum_{p\in\mathcal{I}}\biggl[\frac{\beta_p}{(x-p)^2}+\frac{d_p}{(x-p)}\biggr]
=
\frac{\beta_t}{u^2}+\frac{d_t}{u}
+\sum_{p\in\mathcal{I}\setminus\{t\}}
\biggl[\frac{\beta_p}{(u+t-p)^2}+\frac{d_p}{u+t-p}\biggr]\\
&=\frac{\beta_t}{u^2}+\frac{d_t}{u}
+\sum_{j=0}^\infty(-1)^{j} \sum_{p\in\mathcal{I}\setminus\{t\}}\frac{\beta_p(j+1)+d_p(t-p)}{(t-p)^{j+2}}u^j
=\sum_{j=0}^{\infty}\eta_j u^{j-2},
\end{align*}
}%
where
\[\eta_0=\beta_t,\qquad \eta_1=d_t,\qquad
 \eta_j=(-1)^{j} \sum_{p\in\mathcal{I}\setminus\{t\}}\frac{\beta_p(j-1)+d_p(t-p)}{(t-p)^{j}},\qquad j\geq 2,
 \]
i.e., $\prod_{p\in\mathcal{I}\setminus\{t\}}(t-p)^j\eta_j\in \Q_\theta[\boldsymbol d,\boldsymbol t]$.

Consequently, $y''(x)=Q(x)y(x)$ if and only if
\begin{equation}\label{e0q-14}
j(j-\theta_t)c_j=\sum_{k=0}^{j-1}\eta_{j-k}c_k,\qquad\forall j\geq 0,
\end{equation}
where $(j-\alpha_t/2)(j-\alpha_t/2-1)-\eta_0=j(j-\alpha_t-1)=j(j-\theta_t)$ is used. Clearly, (\ref{e0q-14}) holds automatically for $j=0$, and (\ref{e0q-14}) with $j=1$ leads to $c_1=\frac{d_t}{1-\theta_t}$. By an induction argument, for any $2\leq j\leq \theta_t-1$, $c_j$ can be uniquely solved by (\ref{e0q-14}) as
\[c_j=r_j\bigl[d_{t}^j+\text{l.o.t. in $\boldsymbol d$}\bigr]\in \Q_\theta(\boldsymbol t)[\boldsymbol d]\]
with total degree $j$ in $\boldsymbol d$,
where $r_j\in\Q_\theta\setminus\{0\}$. Furthermore,
\[\prod_{p\in\mathcal{I}\setminus\{t\}}(t-p)^j c_j\in \Q_\theta[\boldsymbol d,\boldsymbol t].\]
Consequently, (\ref{e0q-14}) with $j=\theta_t$ leads to the existence of $r_{\theta_t}\in\Q_\theta\setminus\{0\}$ and a polynomial
\[\hat{P}(\boldsymbol d, \boldsymbol t)=d_t^{\theta_t}+\text{l.o.t.\ in $\boldsymbol d$}\in \Q_\theta(\boldsymbol t)[\boldsymbol d]\]
with total degree $\theta_t$ in $\boldsymbol d$ such that
\[
 \sum_{k=0}^{\theta_t-1}\eta_{j-k}c_k=r_{\theta_t}
 \hat{P}(\boldsymbol d, \boldsymbol t).
\]
Then it is standard by the Frobenius theory that (\ref{00eq: ODE on
 C}) is apparent at the singularity $t$ if and only if
$\hat{P}(\boldsymbol d, \boldsymbol t)=0$.
Furthermore, the above argument also implies the existence of an integer $0\leq N\leq \theta_t$ such that
\[
 P(\boldsymbol d, \boldsymbol t):=\hat{P}(\boldsymbol d, \boldsymbol
 t)\prod_{p\in\mathcal{I}\setminus\{t\}}(t-p)^{N}\in
 \Q_\theta[\boldsymbol d,\boldsymbol t]
\]
and $P(\boldsymbol d, \boldsymbol t)$, $\prod_{p\in\mathcal{I}\setminus\{t\}}(t-p)^{N}$ are coprime.
\end{proof}
